\chapter{Il problema dell'orario scolastico}
\label{ch:problema_timetabling}

\epigraph{``Educational timetabling is one of the most studied
problems in operations research, yet remains one of the
hardest in practice.''}{Andrea Schaerf,
\emph{A survey of automated timetabling}~\parencite{schaerf1999}}

\lettrine[lines=3]{P}{rima} di entrare nel merito di come
$\pi$Tantum costruisce un orario, conviene fermarsi a guardare
con attenzione il problema che si sta affrontando. Lo si pu\`o
descrivere a parole in poche righe, ma chi ci ha provato sa
che quella semplicit\`a \`e ingannevole: si tratta di assegnare a
ogni lezione settimanale uno slot di tempo e un'aula, in modo
che nessun docente, nessuna classe e nessuna aula si trovino
impegnati in due luoghi nello stesso istante, e che il monte
ore di ogni materia in ogni classe sia rispettato. Detto cos\`i
sembra una pratica routinaria, di quelle che si fanno con un
foglio di calcolo in mezza giornata. Nella realt\`a, una scuola
di dimensioni medie produce un'esplosione combinatoria di
possibilit\`a che renderebbe impraticabile qualunque tentativo
di enumerazione esaustiva, e le richieste personali dei
docenti, le incompatibilit\`a delle aule, le esigenze
didattiche delle classi, i vincoli sindacali e le delibere
collegiali si stratificano gli uni sugli altri fino a rendere
il problema, in linea di principio, intrattabile. Vale dunque
la pena dedicare un capitolo intero a comprenderne la natura
prima di affrontarlo, perch\'e una cattiva comprensione del
problema porta inevitabilmente a una cattiva soluzione.

Lo scopo di queste pagine \`e dunque triplice. In primo luogo,
si vuole dare un'idea quantitativa della dimensione del
problema, perch\'e i numeri aiutano a capire perch\'e
``provarci'' non \`e una strategia praticabile. In secondo
luogo, si vogliono distinguere con cura i diversi tipi di
vincolo (rigidi, morbidi, preferenze, obbligazioni), perch\'e
\`e proprio questa distinzione che rende possibile costruire
un sistema utilizzabile. In terzo luogo, si vuole tracciare
una piccola storia dei metodi che la comunit\`a di ricerca
operativa ha sviluppato negli ultimi cinquant'anni, in modo
che il lettore possa collocare $\pi$Tantum in una tradizione
e capire quali idee gli sono state ereditate dai predecessori.

\section{Quante combinazioni possibili}

Si consideri una scuola italiana di dimensione media: trenta
classi, sessanta docenti, sei giorni di lezione, sei ore al
giorno. Ogni classe ha tipicamente una trentina di ore
settimanali, e l'insieme delle lezioni da collocare \`e
dunque dell'ordine di novecento. Per ognuna di queste lezioni
\`e necessario decidere quale materia mettere in quale slot e
quale docente farla insegnare. Anche limitando le scelte ai
soli docenti compatibili (cio\`e quelli abilitati alla materia
in questione e con sufficienti ore residue), per ogni lezione
restano in media una decina di possibilit\`a. Lo spazio delle
soluzioni cresce dunque come dieci elevato a novecento, un
numero che si scrive con un uno seguito da novecento zeri e
che non ha alcun significato fisico immaginabile. Per dare
un termine di paragone, si stima che il numero di atomi
nell'universo osservabile sia attorno a dieci elevato alla
ottantesima potenza: lo spazio delle assegnazioni di una
scuola media \`e quindi $10^{820}$ volte pi\`u grande
dell'universo intero. Esplorarlo per enumerazione cieca non
\`e soltanto impraticabile sui computer di oggi: \`e
fisicamente impossibile.

Questo \`e il motivo per cui non basta ``provare le
combinazioni'', anche con un computer veloce. Servono
algoritmi che, anzich\'e esaminare ogni possibilit\`a, scartino
sotto-spazi enormi senza neppure visitarli, sulla base di
deduzioni logiche elementari. \`E quello che fa la
\emph{propagazione di vincoli}, una tecnica che vedremo nel
dettaglio nel capitolo dedicato a CP-SAT
(cap.~\ref{ch:metodo_cpsat}). In alternativa, si possono
accettare soluzioni \emph{non garantite ottime}, ottenute
attraverso processi di ricerca locale che migliorano una
soluzione iniziale per piccoli passi: \`e l'approccio delle
metaeuristiche, di cui parler\`a il
cap.~\ref{ch:metodo_metaeuristiche}. Si pu\`o anche
spezzare il problema in sotto-problemi pi\`u piccoli e
risolverli quasi indipendentemente, ricucendoli alla fine: \`e
il principio della decomposizione, che $\pi$Tantum applica in
diverse forme (cap.~\ref{ch:metodo_spettrale} per la
decomposizione spettrale, cap.~\ref{ch:metodo_lagrangian} per
quella lagrangiana). Tutte queste tecniche convivono e
collaborano nella pipeline del sistema, che combina ci\`o che
ognuna fa meglio.

\section{Vincoli rigidi e vincoli morbidi}

Una difficolt\`a almeno altrettanto grave della dimensione
combinatoria deriva dalla \emph{natura} dei vincoli che
intervengono. Non sono tutti uguali: alcuni sono assoluti,
non negoziabili, e una soluzione che li violi non \`e neppure
una soluzione, mentre altri esprimono preferenze, desideri o
abitudini di servizio, e si vorrebbero rispettare ma su di
essi si pu\`o transigere se le circostanze lo richiedono.
Ignorare questa distinzione e trattare tutto come ugualmente
inderogabile produce, come si pu\`o immaginare, sistemi
infattibili al primo conflitto fra preferenze, e l'utente si
ritrova davanti a un solver che dichiara ``nessuna soluzione
esistente'' senza la minima indicazione di cosa fare.

In $\pi$Tantum si \`e adottata, fin dall'inizio, una
distinzione canonica della letteratura~\parencite{schaerf1999,carter1998}
fra vincoli rigidi e morbidi, raffinata in cinque livelli
distinti rappresentati nell'interfaccia da una piccola
palette di colori. Il livello \emph{HARD} (rosso) rappresenta
i vincoli inderogabili: i due esempi pi\`u immediati sono
``due classi non possono trovarsi nella stessa aula nello
stesso slot'' e ``un docente non pu\`o tenere due lezioni
contemporaneamente''. Il livello \emph{SOFT} (giallo)
esprime una preferenza negativa: si vorrebbe evitare quella
configurazione (per esempio una sesta ora di matematica al
sabato pomeriggio) ma se serve per chiudere l'orario la si
pu\`o accettare, pagando un costo che entra nella funzione
obiettivo. Il livello \emph{PREFERRED} (blu) \`e simmetrico
rispetto al precedente: indica una preferenza positiva, una
configurazione che si vorrebbe vedere ricorrere (per esempio
il laboratorio di fisica in due ore consecutive) e che riduce
il costo se viene rispettata. Il livello \emph{ENFORCED}
(verde scuro) \`e ancora pi\`u rigido di HARD ma in senso
opposto: si dichiara che quel particolare slot \emph{deve}
essere usato, non \`e ammessa l'alternativa. Infine il livello
\emph{ALLOWED} (verde chiaro) rappresenta la situazione
neutra di default: nessuna preferenza in nessun senso.

Questa stratificazione, che a prima vista pu\`o apparire
sovrabbondante, si rivela quasi sempre necessaria nella
pratica. Una scuola italiana ha decine di vincoli che
appartengono naturalmente al livello SOFT (le preferenze
mattutine dei docenti, la concentrazione delle materie
teoriche nelle prime ore, l'evitamento dei buchi nelle
classi), accanto a quelli inderogabili. Fingere che siano
tutti HARD vorrebbe dire produrre sistemi cronicamente
infattibili; fingere che siano tutti SOFT vorrebbe dire
produrre sistemi che violano regole contrattuali. La
distinzione precisa \`e ci\`o che permette di ottenere
soluzioni che funzionano davvero.

\section{Una breve storia delle idee}

Il problema del timetabling \`e oggetto di studio sistematico
fin dagli anni sessanta del Novecento, quando le scuole pi\`u
grandi cominciarono a cercare metodi sistematici per evitare
i lunghi pomeriggi di settembre passati a comporre l'orario
con scotch e cartoncini ritagliati. I primi tentativi furono
basati su algoritmi di colorazione di grafi, perch\'e una
versione semplificata del problema (assegnare slot in modo
che lezioni in conflitto ricevano slot diversi) si lascia
ricondurre con eleganza al colorare i nodi di un grafo in
modo che nodi adiacenti abbiano colori distinti. Il problema
della colorazione dei grafi \`e stato uno dei primi
classificati come \emph{NP-completo} da Karp nel 1972, e per
riduzione anche il timetabling fu identificato come problema
intrinsecamente difficile. La sigla NP-completo, per chi non
\`e abituato al gergo della complessit\`a computazionale, vuole
dire che non si conosce alcun algoritmo capace di risolvere
\emph{tutte} le istanze del problema in tempo crescente
polinomialmente con la loro dimensione, e che questa
difficolt\`a \`e condivisa con un'intera famiglia di problemi
ritenuti intrattabili in generale~\parencite{garey1979}.

Fortunatamente le istanze \emph{reali} di una scuola italiana
non sono i casi peggiori previsti dalla teoria: la struttura
dei vincoli \`e molto regolare (gli indirizzi sono pochi, le
classi di concorso sono ancora pi\`u poche, le aule
specializzate sono distribuite in modo prevedibile), e questa
regolarit\`a si lascia sfruttare dagli algoritmi moderni in
modo molto efficace. Negli anni ottanta e novanta, alle
tecniche di backtracking si aggiunsero le \emph{metaeuristiche}:
il \emph{simulated annealing} di Kirkpatrick, Gelatt e
Vecchi del 1983~\parencite{kirkpatrick1983}, ispirato alla
metallurgia, e il \emph{tabu search} di Glover del
1986~\parencite{glover1986}, che introdusse l'idea di una
memoria delle mosse recenti per guidare la ricerca. Negli
anni duemila i \emph{constraint solver} acquisirono maturit\`a
industriale: prodotti come GeCode, ECLiPSe e poi CP-SAT di
Google permisero di affrontare istanze prima inarrivabili.
Pisinger e Ropke nel 2006 introdussero la \emph{Adaptive
Large Neighborhood Search}~\parencite{ropke2006}, che combina
distruzioni e ricostruzioni intelligenti su porzioni
significative della soluzione. La conferenza biennale PATAT
(\emph{Practice and Theory of Automated Timetabling}), attiva
dal 1995, divenne il riferimento mondiale per i ricercatori
del settore~\parencite{patat}.

$\pi$Tantum si colloca in questa tradizione. Non inventa
metodi nuovi: combina con cura quelli esistenti in una
pipeline orchestrata che cerca di sfruttare i punti di forza
di ciascuno. CP-SAT viene usato come motore principale per
la fase di soddisfacimento dei vincoli rigidi, sia nella
fase di matching docente--classe sia nella collocazione
settimanale. La decomposizione spettrale, ispirata alle
intuizioni di Fiedler del 1973~\parencite{fiedler1973}, viene
usata per spezzare le scuole pi\`u grandi in cluster
naturali, in modo che CP-SAT possa lavorare su sotto-problemi
trattabili. Le metaeuristiche raffinano la soluzione
ottenuta, riducendo le violazioni dei vincoli morbidi.
Il pre-controllo di Hall, basato su un teorema del 1935 che
sembra antico ma che resta sorprendentemente attuale,
verifica in pochi millisecondi se il problema \`e almeno
formalmente fattibile, evitando di lanciare il solver su
istanze infeasibili dall'inizio.

\section{Le quattro famiglie di approcci}

La letteratura sul timetabling distingue tradizionalmente
quattro famiglie di approcci, ciascuna con i propri vantaggi
e i propri limiti. La prima \`e la \emph{programmazione a
vincoli}, che modella il problema come insieme di variabili
con domini finiti e di vincoli che le legano, lasciando poi
che un solver come CP-SAT si occupi di trovare un'assegnazione
che li soddisfi tutti. Il vantaggio di questo approccio \`e la
sua estrema flessibilit\`a: aggiungere un vincolo nuovo
significa aggiungere una riga di codice e nessuna riprogettazione
strutturale del modello. Lo svantaggio \`e che la scalabilit\`a
diventa critica oltre la soglia di una settantina di classi:
il numero di variabili boolean equivalenti che il solver deve
gestire cresce e i tempi di risoluzione esplodono.

La seconda famiglia \`e la \emph{programmazione lineare intera}
(MILP, \emph{Mixed Integer Linear Programming}), che formula
il problema come sistema di disuguaglianze lineari con
variabili binarie e lo risolve con prodotti commerciali come
Gurobi o CPLEX, oppure con alternative open-source come HiGHS.
Il vantaggio teorico \`e che si possono ottenere certificati
di ottimalit\`a, almeno su istanze piccole. Lo svantaggio
pratico \`e che la formulazione MILP del timetabling tende a
richiedere milioni di variabili binarie e diventa pesante.
$\pi$Tantum non utilizza MILP nella pipeline principale ma
prevede uno scheletro di generazione di colonne (\emph{column
generation}) che si appoggia a un master LP risolto con HiGHS,
applicabile alle istanze pi\`u grandi (cap.~\ref{ch:metodo_lagrangian}).

La terza famiglia \`e quella delle \emph{metaeuristiche}.
Si parte da una soluzione iniziale, ottenuta con un metodo
qualunque (anche un greedy molto rudimentale), e la si
migliora iterativamente attraverso operatori di ricerca
locale: scambi di lezioni, perturbazioni casuali, accettazioni
parziali di mosse peggioranti. Il vantaggio \`e la scalabilit\`a:
le metaeuristiche restituiscono \emph{sempre} una soluzione,
anche su istanze enormi, e in tempi controllabili. Lo
svantaggio \`e che non si ha alcuna garanzia di ottimalit\`a:
la soluzione finale \`e di solito buona, ma raramente la
migliore possibile.

La quarta famiglia \`e quella delle \emph{decomposizioni}.
Invece di affrontare il problema monolitico, lo si spezza in
sotto-problemi pi\`u piccoli, lo si risolve indipendentemente,
e lo si ricuce alla fine. La decomposizione pu\`o essere
\emph{spettrale}, basata sulla struttura algebrica del grafo
bipartito che lega classi e docenti; \emph{lagrangiana},
basata sul rilassamento dei vincoli che legano i sotto-problemi
e sulla loro penalizzazione progressiva nell'obiettivo;
\emph{di Dantzig-Wolfe}, basata sulla generazione iterativa
di nuove variabili nel master problem. $\pi$Tantum implementa
tutte e tre, con preferenze diverse a seconda della
dimensione dell'istanza.

\section{La filosofia di $\pi$Tantum}

La scelta progettuale fondamentale di $\pi$Tantum \`e di
\emph{non scegliere} fra le quattro famiglie sopra descritte,
ma di combinarle in una pipeline a quattro fasi orchestrate.
Ogni fase utilizza il metodo pi\`u appropriato al suo compito
e produce un risultato che la fase successiva pu\`o usare
come punto di partenza. Si tratta di un approccio non
particolarmente originale dal punto di vista della ricerca
(\`e ci\`o che molti sistemi commerciali e accademici fanno),
ma di un approccio efficace, e che ha il pregio di rendere il
sistema modulare: si pu\`o sostituire una fase con
un'implementazione alternativa senza dover toccare le altre.

La prima fase, denominata \emph{Phase A} nella terminologia
del progetto, si occupa dell'\emph{assegnamento}: per ogni
classe e per ogni materia, decidere quale docente la insegner\`a
e quante ore. \`E un problema di matching nel grafo bipartito
classi--docenti, vincolato dalle compatibilit\`a fra classi di
concorso e materie, dai monte ore disponibili per ciascun
docente, e dalle indisponibilit\`a per giorno. Si risolve con
un algoritmo greedy seguito da un passo CP-SAT che ottimizza
la distribuzione. Prima di lanciare la fase, $\pi$Tantum
esegue il \emph{pre-controllo di Hall}, che verifica
in pochi millisecondi se il problema \`e almeno formalmente
fattibile dal punto di vista dei matching: se non lo \`e, il
sistema lo segnala all'utente con l'indicazione del
sottoinsieme problematico, evitando un'attesa inutile su
CP-SAT.

La seconda fase, \emph{Phase B}, si occupa del
\emph{timetabling vero e proprio}: collocare ogni lezione in
uno slot della settimana. Per le scuole grandi, prima di
lanciare CP-SAT si applica la \emph{decomposizione spettrale},
che spezza il problema in cluster relativamente indipendenti.
CP-SAT risolve allora ogni cluster in poche decine di
secondi, e un passo finale di ricucitura armonizza le
interfacce fra cluster.

La terza fase, \emph{Phase C}, si occupa del
\emph{raffinamento}: ALNS, simulated annealing, tabu search e
iterated local search lavorano sulla soluzione di Phase B per
ridurre ulteriormente il valore della funzione obiettivo
SOFT. Si tratta di una fase opzionale ma quasi sempre
attivata, che riduce il costo SOFT di un dieci-quindici per
cento in pochi minuti aggiuntivi.

La quarta fase, \emph{Phase D}, si occupa
dell'\emph{assegnazione delle aule}: dato che ogni lezione
ha gi\`a uno slot e una classe, resta da decidere in quale
aula si terr\`a, rispettando i tipi (lab, palestra, aula
normale), le capienze, le preferenze classe-aula e le
compatibilit\`a materia-aula. Si tratta di un problema CP-SAT
relativamente piccolo, che si risolve in pochi secondi.

L'intera pipeline, dal momento in cui l'utente preme il
bottone ``Pipeline completa'' al momento in cui il
sistema mostra l'orario finito, dura tipicamente un minuto su
una scuola piccola, due minuti su una media, dieci minuti su
una grande, e fino a un quarto d'ora sulle pi\`u grandi
testate. Sono numeri che andranno verificati nel
cap.~\ref{ch:benchmarks}, dove riporto i risultati
sperimentali ottenuti su una macchina consumer con i cinque
profili di test.

\section{Le peculiarit\`a del caso italiano}

Vale la pena chiudere questo capitolo con un'osservazione che
torner\`a spesso nelle pagine successive: i benchmark
internazionali del timetabling, codificati in formati standard
come ITC-2007, ITC-2011 e XHSTT~\parencite{mccollum2010},
sono importanti per il confronto fra ricercatori, ma non
catturano alcune peculiarit\`a del sistema scolastico italiano
che invece $\pi$Tantum deve modellare con precisione. La
prima peculiarit\`a sono gli \emph{indirizzi di studio} (Liceo
Classico, Liceo Scientifico, Linguistico, Istituto Tecnico
nelle sue varie articolazioni, Istituti Professionali), che
hanno monte ore diversi anche per la stessa materia. Una
classe ``3A scientifico'' fa matematica con un orario diverso
da una ``3A classico'' o da una ``3A ITIS Informatica'',
nonostante che il nome della materia sia lo stesso. La seconda
peculiarit\`a sono le \emph{classi di concorso}, codificate
con sigle alfanumeriche (A027, A019, eccetera), che limitano
quali docenti possono insegnare quali materie e in quali
indirizzi. Un docente di A027 pu\`o insegnare matematica e
fisica al biennio scientifico ma non lettere al classico, e
queste compatibilit\`a sono codificate in un grande database
ministeriale che $\pi$Tantum sintetizza nelle sue tabelle.

La terza peculiarit\`a \`e la presenza diffusa di
\emph{compresenze}: in molti indirizzi tecnici, le ore di
laboratorio richiedono che due docenti diversi siano nella
stessa classe nello stesso slot (l'insegnante teorico e
l'insegnante tecnico-pratico). La compresenza \`e un vincolo
HARD non banale, perch\'e non si limita a chiedere ``due
docenti'' ma chiede ``proprio quei due docenti, contemporaneamente''.
La quarta peculiarit\`a sono i \emph{gruppi articolati}: un
sottoinsieme di studenti, eventualmente proveniente da pi\`u
classi di base, che si riuniscono in slot dedicati per
attivit\`a specifiche (la seconda lingua straniera del
linguistico, la lezione di religione cattolica con
l'alternativa per chi non si avvale, i corsi di recupero in
itinere). La quinta peculiarit\`a sono i \emph{vincoli
sindacali}: il giorno libero contrattuale CCNL, il monte ore
massimo settimanale, la priorit\`a basata sul punteggio di
graduatoria. Sono vincoli HARD non negoziabili, ai quali si
aggiungono i vincoli logici espressi via DSL dall'utente.

Tutte queste peculiarit\`a sono codificate come entit\`a o
vincoli di prima classe nel modello dati di $\pi$Tantum
(cap.~\ref{ch:modello_dati}), che \`e stato progettato a
partire da esse e non come adattamento di un modello
accademico generico. \`E forse il contributo pi\`u distintivo
del progetto: non un solver pi\`u veloce, ma un solver che
parla la lingua del sistema scolastico italiano.

Detto questo, abbiamo dato uno sguardo d'insieme al problema.
Nel prossimo capitolo descriveremo $\pi$Tantum dall'alto:
quali sono i suoi tre layer (solver, web UI, prototipi
legacy), come comunicano fra loro, e quale flusso di lavoro
quotidiano un coordinatore d'orario segue per usarlo.
Affronteremo poi, nei capitoli successivi, l'orario visto
come oggetto strutturale (le entit\`a e le relazioni che
codifica), l'interfaccia utente, i casi d'uso ricorrenti, e
infine ciascuno dei metodi algoritmici che fanno funzionare
il sistema.
